\chapter{Weil positivity as the Thorin condition}\label{chap:wt}\chaptermark{Weil positivity and Thorin}

Weil's criterion states RH as the positivity of a functional built from the explicit formula. On Gaussian test functions and their derivatives in the heat time this functional is the heat trace $W(t)$ of the zeros, and positivity becomes the Thorin condition of the reciprocal $\Xi$ law: RH holds if and only if $W$ is completely monotone. The explicit formula writes $W=\mathcal A-\mathcal P$ for every $t>0$, with an archimedean part $\mathcal A$ and a prime part $\mathcal P$, and both have an extreme value form: $\mathcal A$ is $u$ times the L\'evy measure of a Gumbel type variable of order $\frac14$, and $\mathcal P$ is a positive weighted series of Fr\'echet distribution functions, one for each prime power.

The chapter is based on \cite{P26wt}. Its numbers were recomputed from the mathematics alone, without the code of the source. Every number printed below is one of those recomputations unless it is marked unchecked, and where the recomputation differed from the source the recomputed value is printed. Statements of papers that were not read are marked unchecked.

Notation. The heat trace of the zeros is the $W(t)$ of Section~\ref{ch:clock}; the $\Pi+\mathcal A$ of Table~\ref{tab:ch14:W} is written $\mathcal A(t)$ here, and $\mathcal P(t)$ is the prime part. The Gaussian kernel is $g_t(u)=(4\pi t)^{-1/2}e^{-u^2/4t}$. Chebyshev's function is $\psi_{\mathrm C}(y)=\sum_{n\le y}\Lambda(n)$, and $\psi=\Gamma'/\Gamma$. The normalised prime error is $D(u)=(\psi_{\mathrm C}(e^u)-e^u)e^{-u/2}$. The shifted functional of Section~\ref{wt:sec:sin2} is $V(t,a)$. Zeros are indexed as in Section~\ref{ch:clock} (one zero of each pair, with multiplicity), with local symbols: $\rho_j=\beta_j+i\tau_j$ with $\tau_j>0$ (the $\gamma_k$ of Section~\ref{ch:clock}); $\gamma_j=\tau_j+i(\frac12-\beta_j)$ is complex, so that $\rho_j=\frac12+i\gamma_j$ and $\gamma_j$ is real exactly when $\rho_j$ is on the critical line; and $\lambda_j=\gamma_j^2$ is the $a_k=-(\rho_k-\frac12)^2$ of the book. A sum over positive real $\gamma$ is written only under RH. The first window root is $r_*$ with $\lambda_*=r_*^2$.

The explicit formula \eqref{eq:ch2:Warith} is also a heat equation. With $p_x(u)=e^{-u^2/4x}/(2\sqrt{\pi x})$, $W(x)=e^{x/4}+\langle\mathcal D,p_x\rangle$, where $\mathcal D=-\frac12\log\pi\,\delta_0+\frac12\mathcal M-\frac12\sum_{n\ge2}\Lambda(n)n^{-1/2}(\delta_{\log n}+\delta_{-\log n})$ is an even distribution and $\mathcal M$ is the tempered distribution whose Fourier transform is $\re\psi(\frac14+\frac{ir}2)$, normalised so that $\langle\mathcal M,p_x\rangle=\frac1{2\pi}\int e^{-r^2x}\re\psi(\frac14+\frac{ir}2)dr$. So $W(x)-e^{x/4}$ is the solution at time $x$ and position $0$ of $\partial_xu=\partial_u^2u$ started from the prime distribution $\mathcal D$ (proved; the label is Kolmogorov's backward equation, and it is a rewriting of the formula). On the side of Kolmogorov's theta law the initial condition is a point mass and the walls absorb, while here it is the prime distribution on the whole line (an unrefereed note of Polson, November 2026).

Results already in the book are cited, not restated: the heat trace, its Frullani and Weil forms and the equivalence of RH with complete monotonicity (Section~\ref{ch:clock}, Proposition~\ref{prop:rigid}, Equation~\eqref{eq:ch2:Warith}), the digits of Tables~\ref{tab:ch14:W} and~\ref{tab:ch2:W}, the theorems that $W>0$ and $W'<0$ for every $t>0$ (Theorem~\ref{thm:ch12:idsd}), that $W$ is completely monotone on $(0,1.29]$ with no information on the zeros below $0.019$ (Theorem~\ref{thm:ch12:cmint}) and that $(-1)^kW^{(k)}>0$ for all $t>0$ and $k\le1.16\times10^{25}$ (Theorem~\ref{thm:ch12:window}), the tent version of the sine squared lemma (Theorem~\ref{yo:thm:D}), the Gumbel kernels of fractional order (Section~\ref{ev:sec:decomp}) and the Prony extraction of the atoms (Section~\ref{sec:ch18:reach}). The results of the chapter are these.
\begin{enumerate}
\item RH holds if and only if $W$ is completely monotone, if and only if the Thorin measure $\sum_j\delta_{\lambda_j}$ is supported on $(0,\infty)$, if and only if $\mathcal A-\mathcal P$ is completely monotone (Theorem~\ref{wt:thm:main}, proved). For the converse, the Laplace transform of $W$ is meromorphic with a nonzero residue at every $-\lambda_j$, and a nonreal $\lambda_j$ cannot be the pole of a Stieltjes transform.
\item The archimedean density is $u$ times the L\'evy density of $-\frac12\log G_{1/4}$, and the prime part is a weighted series of Fr\'echet distribution functions with scales $(\log n)^2/4$ (Section~\ref{wt:sec:gumbel}, proved).
\item $W$ is computed to 110 digits from $\Gamma$ and the primes; Prony's method on 16 such samples gives $\gamma_1,\gamma_2,\gamma_3$ to more than 30 digits and the eighth zero to 4 digits (Table~\ref{wt:tab:primezeros}, computed).
\item The archimedean part $\mathcal A$ has a spectral density with exactly one sign change, at $\lambda_*=39.562\approx4\pi^2$, where $r_*$ is the critical point of the Riemann Siegel theta function. The first sign change $t_n$ of $(-1)^n\mathcal A^{(n)}$ satisfies $t_n=n/\lambda_*+\beta+\delta/n+O(n^{-2})$ with explicit $\beta$ and $\delta$ (Propositions~\ref{wt:prop:window} and~\ref{wt:prop:expansion}, proved; the table is computed). Complete monotonicity of order $n$ of $\mathcal A$ holds only on $(0,t_1]$.
\item Each prime power added extends the order one window to about $\log q/2\gamma_1$ (computed, a heuristic balance). A one sided certificate from the primes up to $N_0$ is proved (Proposition~\ref{wt:prop:cert}); the windows tabulated from it are computed in floating point on a grid, not certified.
\item The variance of $\psi_{\mathrm C}$ in short intervals follows the GUE prediction and not the Poisson level (computed, with a heuristic prediction). RH is equivalent to finite discounted energy of $D$ (Proposition~\ref{wt:prop:energy}, proved); the mean square constant needs simple zeros.
\item The first two cumulants of the reciprocal $\Xi$ law are read off the primes, and RH is equivalent to the positive semidefiniteness of two Hankel families of the normalised cumulants (Section~\ref{wt:sec:cumulants}, proved).
\item RH holds if and only if $V(t,a)\ge0$ for every $a\ge0$ at one fixed heat time $t$ (Theorem~\ref{wt:thm:sin2}, proved by Landau's theorem for nonnegative functions).
\end{enumerate}
The open problems close the chapter.
